% The 23rd Huawei Cup China Post-Graduate Mathematical Contest in Modeling
% LaTeX template converted from the supplied Word paper template.
% Compile with XeLaTeX or LuaLaTeX.

\documentclass[11pt,a4paper]{article}

\usepackage[a4paper,top=3.0cm,bottom=1.75cm,left=2.25cm,right=2.25cm,footskip=1.9cm]{geometry}
\usepackage{iftex}
\ifPDFTeX
  \usepackage[T1]{fontenc}
  \usepackage[utf8]{inputenc}
  \usepackage{lmodern}
  \newcommand{\stxinwei}{\sffamily}
  \newcommand{\lishu}{\rmfamily}
  \newcommand{\tnr}{\rmfamily}
\else
  \usepackage{fontspec}
  \usepackage{xeCJK}
  \setmainfont{Times New Roman}
  \setCJKmainfont{SimSun}
  \newfontfamily\stxinwei{STXinwei}
  \newfontfamily\lishu{LiSu}
  \newfontfamily\tnr{Times New Roman}
\fi
\usepackage{amsmath,amssymb,mathtools}
\usepackage{graphicx}
\usepackage{xcolor}
\usepackage{booktabs}
\usepackage{tabularx}
\usepackage{array}
\usepackage{multirow}
\usepackage{enumitem}
\usepackage{microtype}
\usepackage{caption}
\usepackage{subcaption}
\usepackage{float}
\usepackage{fancyhdr}
\usepackage{titlesec}
\usepackage[hidelinks]{hyperref}

\definecolor{HuaweiRed}{HTML}{D71920}
\definecolor{ContestGreen}{HTML}{2B9B2B}
\definecolor{RuleGray}{HTML}{777777}

\setlength{\parindent}{2em}
\setlength{\parskip}{0pt}
\setlength{\emergencystretch}{2em}
\renewcommand{\baselinestretch}{1.08}

% Four numbered heading levels: section, subsection, subsubsection, paragraph.
\setcounter{secnumdepth}{4}
\setcounter{tocdepth}{4}
\captionsetup{font=small,labelfont=bf,labelsep=period}

\titleformat{\section}
  {\normalfont\large\bfseries}{\thesection}{0.7em}{}
\titleformat{\subsection}
  {\normalfont\normalsize\bfseries}{\thesubsection}{0.6em}{}
\titleformat{\subsubsection}
  {\normalfont\normalsize\bfseries}{\thesubsubsection}{0.6em}{}
\titleformat{\paragraph}
  {\normalfont\small\bfseries}{\theparagraph}{0.6em}{}
\titlespacing*{\section}{0pt}{1.1em}{0.45em}
\titlespacing*{\subsection}{0pt}{0.8em}{0.3em}
\titlespacing*{\subsubsection}{0pt}{0.65em}{0.25em}
\titlespacing*{\paragraph}{0pt}{0.55em}{0.2em}

\pagestyle{fancy}
\fancyhf{}
\fancyfoot[C]{\thepage}
\renewcommand{\headrulewidth}{0pt}
\renewcommand{\footrulewidth}{0pt}

\fancypagestyle{cover}{%
  \fancyhf{}
  \fancyfoot[C]{%
    {\stxinwei\fontsize{12}{14}\selectfont
      China Postgraduate Innovation \& Practice Competitions\\[-0.1em]}
    {\normalfont\fontsize{9}{11}\selectfont 0}%
  }
  \renewcommand{\headrulewidth}{0pt}
  \renewcommand{\footrulewidth}{0pt}
}

% ---------- Editable paper metadata ----------
\newcommand{\University}{Your University}
\newcommand{\TeamNumber}{Your Team Number}
\newcommand{\TeamMemberOne}{First Author}
\newcommand{\TeamMemberTwo}{Second Author}
\newcommand{\TeamMemberThree}{Third Author}
\newcommand{\PaperTitle}{Your Paper Title}
\newcommand{\AbstractText}{}
\newcommand{\Keywords}{}

\newcommand{\ContestTitle}{%
  {\stxinwei\fontsize{22}{27}\selectfont\bfseries
    The 23rd “Huawei Cup” China Post-Graduate\\[-0.05em]
    Mathematical Contest in Modeling}%
}

\newcommand{\CoverInfoTable}{%
  {\normalfont\fontsize{18}{25}\selectfont\bfseries
   \setlength{\arrayrulewidth}{0.75pt}
   \setlength{\tabcolsep}{0pt}
   \renewcommand{\arraystretch}{1.15}
   \begin{tabular}{@{}p{3.25cm}p{11.25cm}@{}}
     University & \hrulefill\\ \hline
     Team & \hrulefill\\
     Number & \hrulefill\\ \hline
     \multirow{3}{3.25cm}{\centering Team\\Members} & 1.\ \hrulefill\\
       & 2.\ \hrulefill\\
       & 3.\ \hrulefill\\ \hline
   \end{tabular}}
}

\newcommand{\DocCoverPage}{%
  \begin{center}
    \begin{tabular}{@{}c@{\hspace{0.5cm}}c@{\hspace{0.5cm}}c@{\hspace{0.5cm}}c@{}}
      
\includegraphics[width=102pt,height=65pt]{figures/cpipc-logo.png} &
      
\includegraphics[width=78.21819pt,height=75.27273pt]{figures/contest-logo.png} &
      
\includegraphics[width=79.44pt,height=79.44pt]{figures/huawei-logo.png} &
      
\includegraphics[width=68.16pt,height=64.32pt]{figures/university-logo.png}\\
    \end{tabular}\par
    \vspace{1.65cm}
    {\stxinwei\fontsize{18}{22}\selectfont\bfseries
      China Postgraduate Innovation \& Practice Competitions\\[0.25em]}
    \ContestTitle
    \vspace{1.55cm}
    \CoverInfoTable
  \end{center}
}

\newcommand{\DocAbstractPage}{%
  \vspace{0.2cm}
  \begin{center}
    \ContestTitle
  \end{center}
  \vspace{1.45cm}
  {\lishu\fontsize{18}{22}\selectfont
    {\tnr Title}{\lishu ：}\hspace{0.8cm}\rule{8.7cm}{0.4pt}\\[0.2cm]
    \begin{center}{\tnr Abstract}{\lishu ：}\end{center}
    \vspace{3.2cm}
    \noindent {\tnr Keywords}{\lishu ：}\par
  }
  \ifx\AbstractText\empty\else
    \vspace{0.4em}\noindent\AbstractText\par
  \fi
  \ifx\Keywords\empty\else
    \noindent\Keywords\par
  \fi
}

\begin{document}

\pagenumbering{arabic}
\setcounter{page}{0}
\thispagestyle{cover}
\DocCoverPage
\clearpage
\DocAbstractPage
\clearpage

% ---------- Table of contents ----------
\pagenumbering{roman}
\tableofcontents
\clearpage
\pagenumbering{arabic}

% ---------- Main text ----------
\section{Introduction}
\subsection{Background and Motivation}
Describe the real-world background and explain why the problem is important.

\subsection{Objectives and Contributions}
State the research objectives and summarize the main contributions in a concise list.

\subsubsection{Research Questions}
List the questions that the model is designed to answer.

\paragraph{Contribution Statement}
Explain the specific methodological or practical contribution in one focused paragraph.

\section{Problem Restatement and Data Description}
\subsection{Problem Restatement}
Restate each subproblem in precise mathematical language without copying the original
question verbatim.

\subsection{Data Sources and Preprocessing}
Describe the data source, sample size, units, missing-value treatment, normalization,
and any other preprocessing steps.

\subsubsection{Data Quality Checks}
Report outlier detection, consistency checks, and the criteria used to retain or remove
observations.

\paragraph{Notation}
Define the symbols used throughout the paper. For example, $x_i$ denotes the value of
the $i$-th decision variable and $\mathcal{X}$ denotes the feasible set.

\section{Assumptions and Model Formulation}
\subsection{Modeling Assumptions}
List assumptions and explain why each is reasonable for the problem.

\subsection{Variables and Parameters}
Define all decision variables, state variables, parameters, and units before using them.

\subsubsection{Objective Function}
Present the optimization objective and explain its interpretation.

\begin{equation}
  \min_{x \in \mathcal{X}}\; f(x)
  \label{eq:objective}
\end{equation}

\paragraph{Constraints}
Give the model constraints and explain the meaning of their bounds and domains.

\section{Solution Method}
\subsection{Algorithm Overview}
Describe the overall solution procedure and include a flowchart when it improves clarity.

\subsection{Detailed Algorithm}
Give sufficient detail for the method to be reproduced, including initialization,
termination criteria, parameter settings, and computational complexity when relevant.

\subsubsection{Numerical Implementation}
Record the software environment, solver, random seed, and important implementation choices.

\paragraph{Reproducibility Note}
Provide the command, configuration, or repository information needed to reproduce the results.

\section{Results and Discussion}
\subsection{Baseline Results}
Present the main results and compare them with a baseline or benchmark.

\subsection{Figures and Figure Captions}
Every figure must be numbered, have an informative caption below it, and be cited in the
main text. Use a caption that explains the plotted quantity, units, and experimental setting.
For example, Fig.~\ref{fig:overall-framework} shows the overall workflow.

\begin{figure}[H]
  \centering
  % Replace the placeholder with, for example:
  % \includegraphics[width=0.82\linewidth]{figures/flow_overall_model.png}
  \fbox{\rule{0pt}{4cm}\rule{0.82\linewidth}{0pt}}
  \caption{Overall framework of the proposed modeling and solution procedure.}
  \label{fig:overall-framework}
\end{figure}

\subsection{Tables and Table Captions}
Table captions are placed above tables in this template. Each table must be numbered,
have clear column headings and units, and be cited in the main text.

\begin{table}[H]
  \centering
  \caption{Definition of the main variables and parameters}
  \label{tab:notation}
  \begin{tabularx}{0.9\textwidth}{@{}lXl@{}}
    \toprule
    Symbol & Description & Unit \\
    \midrule
    $x_i$ & Decision variable for item $i$ & -- \\
    $c_i$ & Cost or coefficient associated with item $i$ & unit \\
    $\mathcal{X}$ & Feasible solution set & -- \\
    \bottomrule
  \end{tabularx}
\end{table}

\subsection{Validation and Sensitivity Analysis}
Evaluate the model using independent data, alternative parameter settings, cross-validation,
or other appropriate tests. Explain which conclusions are robust.

\subsubsection{Sensitivity to Key Parameters}
Show how the results change when the most influential parameters are varied.

\paragraph{Error and Robustness Analysis}
Discuss uncertainty, error sources, limitations, and the conditions under which the model may fail.

\section{Conclusions}
Summarize the main findings, answer the research questions, and state practical implications
and possible future improvements.

% ---------- References ----------
\clearpage
\phantomsection
\section*{References}
\addcontentsline{toc}{section}{References}
Use author--year or numbered citations consistently. Every reference listed below should
be cited in the text, for example, \cite{example-paper}. Replace the sample entries with
verified sources before submission.

\begin{thebibliography}{99}
  \bibitem{example-paper}
  Author, A. A. (Year).
  Title of the article.
  \textit{Journal Name}, volume(issue), pages.

  \bibitem{example-book}
  Author, B. B. (Year).
  \textit{Title of the Book}.
  Publisher.
\end{thebibliography}

% ---------- Optional appendix ----------
\appendix
\section{Supplementary Derivations}
Place lengthy derivations, additional proofs, parameter tables, or supplementary results here.

\section{Additional Figures and Tables}
Place supporting figures and tables here when they are not essential to the main argument.

\end{document}
